% Part: computability
% Chapter: recursive-functions
% Section: pr-functions-computable

\documentclass[../../../include/open-logic-section]{subfiles}

\begin{document}

\olfileid{cmp}{rec}{cmp}
\olsection{بنیادی بازگشتی فنکشن قابلِ حساب نیں}

فرض کرو فنکشن $h$ نوں بنیادی بازگشت راہیں
متعین کیتا گیا اے:
\begin{eqnarray*}
h(\vec x, 0)     & = & f(\vec x) \\
h(\vec x, y + 1) & = & g(\vec x, y, h(\vec x, y))
\end{eqnarray*}
تے فرض کرو فنکشن $f$ تے $g$ قابلِ حساب
نیں۔ (ایتھے اسی $\vec x$ نوں $x_0$،
\dots، $x_{k-1}$ دا مختصر روپ ورتدے
آں۔) تاں $h(\vec x, 0)$ صاف طور اُتے
حساب کیتا جا سکدا اے، کیونکہ ایہ بس
$f(\vec x)$ اے جیہڑا فرض مطابق قابلِ
حساب اے۔ فیر $h(\vec x, 1)$ وی حساب
کیتا جا سکدا اے، کیونکہ $1 = 0 + 1$
تے ایس لئی $h(\vec x, 1)$ بس ایہ اے:
\begin{align*}
 h(\vec x, 1) & = g(\vec x, 0, h(\vec x, 0)) =  g(\vec x, 0, f(\vec x)).
\intertext{اسی ایس طرح اگے ودھ کے حساب کر سکدے آں}
h(\vec x, 2) & = g(\vec x, 1, h(\vec x, 1)) = g(\vec x, 1, g(\vec x, 0, f(\vec x)))\\
h(\vec x, 3) & = g(\vec x, 2, h(\vec x, 2)) = g(\vec x, 2, g(\vec x, 1, g(\vec x, 0, f(\vec x))))\\
h(\vec x, 4) & = g(\vec x, 3, h(\vec x, 3)) = g(\vec x, 3, g(\vec x, 2, g(\vec x, 1, g(\vec x, 0, f(\vec x)))))\\
& \vdots
\end{align*}
لہٰذا عام طور اُتے $h(\vec x, y)$
حساب کرن لئی پہلے $h(\vec x, 0)$،
فیر $h(\vec x, 1)$، \dots، باری باری
حساب کردے جاواں، جد تک $h(\vec x, y)$
تک نہ پہنچ جاواں۔

ایس لئی جے فنکشن $f$ تے $g$ قابلِ
حساب ہون، تاں بنیادی بازگشتی تعریف
اک نواں قابلِ حساب فنکشن دیندی اے۔
اِس طرح جے فنکشن $f$ تے $g_i$ قابلِ
حساب ہون، تاں اوہناں دی ترکیب نال
ملن والا فنکشن وی قابلِ حساب ہوندا
اے۔

چونکہ بنیادی فنکشن $\Zero$، $\Succ$
تے $\Proj{n}{i}$ قابلِ حساب نیں،
تے ترکیب تے بنیادی بازگشت قابلِ حساب
فنکشناں توں قابلِ حساب فنکشن دیندیاں
نیں، اِس لئی ہر بنیادی بازگشتی فنکشن
قابلِ حساب اے۔

\end{document}
